%==============================================================================
% Oefeningen bij Hoofdstuk 1 -- Energie en de eerste hoofdwet
%==============================================================================

\section*{Oefeningen bij Hoofdstuk 1}

\addcontentsline{toc}{section}{Oefeningen bij Hoofdstuk 1}

\begin{enumerate}

\item \textbf{Compressorarbeid (isotherm).} \\
Een industri\"ele compressor comprimeert \SI{100}{\mole} stikstof (ideaal gas) isotherm bij \SI{300}{\kelvin} van \SI{1}{\bar} naar \SI{50}{\bar}. Bereken
\begin{enumerate}
\item[(a)] de benodigde reversibele arbeid,
\item[(b)] de warmte die aan de omgeving moet worden afgevoerd,
\item[(c)] het volume van het gas voor en na de compressie.
\end{enumerate}
\emph{Hint:} voor een isotherm proces geldt $\Delta U = 0$ voor een ideaal gas.

\item \textbf{Compressorarbeid (adiabatisch).} \\
Dezelfde compressor werkt nu adiabatisch ($Q = 0$) en reversibel. De begin temperatuur bedraagt \SI{300}{\kelvin}. Neem $C_{V,m} = \frac{5}{2}R$ voor $\mathrm{N_2}$ (tweatomig ideaal gas).
\begin{enumerate}
\item[(a)] Bereken de eindtemperatuur bij een compressie van \SI{1}{\bar} naar \SI{50}{\bar}.
\item[(b)] Bereken de benodigde arbeid per mol.
\item[(c)] Vergelijk met de isotherme arbeid uit Oefening 1. Verklaar waarom de adiabatische arbeid groter is.
\end{enumerate}
\emph{Hint:} voor een reversibele adiabatische compressie van een ideaal geldt $T_2/T_1 = (p_2/p_1)^{R/C_{p,m}}$.

\item \textbf{Opwarmen van water in een warmtewisselaar.} \\
Een processtroom van \SI{2.0}{\kilogram\per\second} water wordt in een warmtewisselaar opgewarmd van \SI{25}{\celsius} naar \SI{85}{\celsius}. De warmtecapaciteit van water bedraagt $C_p \approx \SI{4.18}{\kilo\joule\per\kilogram\per\kelvin}$ en mag constant worden verondersteld.
\begin{enumerate}
\item[(a)] Bereken de benodigde warmtestroom (in \si{\kilo\watt}).
\item[(b)] Bereken de verandering van de enthalpie van de waterstroom per seconde.
\item[(c)] Waarom mag de enthalpieverandering hier gelijk worden gesteld aan de uitgewisseld warmte?
\end{enumerate}

\item \textbf{Verbrandingsenthalpie via de wet van Hess.} \\
Bereken de standaardreactie-enthalpie van de verbranding van vloeibaar ethanol:
\[
\mathrm{C_2H_5OH(l) + 3\,O_2(g) \rightarrow 2\,CO_2(g) + 3\,H_2O(l)}.
\]
Gegeven zijn de standaardvormingsenthalpie\"en bij \SI{298}{\kelvin}:
\begin{itemize}
\item $\Delta H_f^\circ(\mathrm{C_2H_5OH,l}) = \SI{-277.7}{\kilo\joule\per\mole}$
\item $\Delta H_f^\circ(\mathrm{CO_2,g}) = \SI{-393.5}{\kilo\joule\per\mole}$
\item $\Delta H_f^\circ(\mathrm{H_2O,l}) = \SI{-285.8}{\kilo\joule\per\mole}$
\item $\Delta H_f^\circ(\mathrm{O_2,g}) = 0$
\end{itemize}
Hoeveel warmte komt vrij bij de verbranding van \SI{1}{\kilogram} ethanol?

\item \textbf{Wet van Kirchhoff: reactie-enthalpie bij hoge temperatuur.} \\
De standaardreactie-enthalpie van de ammoniaksynthese
\[
\mathrm{N_2(g) + 3\,H_2(g) \rightarrow 2\,NH_3(g)}
\]
bedraagt $\Delta_r H^\circ(\SI{298}{\kelvin}) = \SI{-92.4}{\kilo\joule\per\mole}$. De molair warmtecapaciteiten (constant verondersteld) zijn:
\begin{itemize}
\item $C_{p,m}(\mathrm{N_2}) = \SI{29.1}{\joule\per\mole\per\kelvin}$
\item $C_{p,m}(\mathrm{H_2}) = \SI{28.8}{\joule\per\mole\per\kelvin}$
\item $C_{p,m}(\mathrm{NH_3}) = \SI{35.1}{\joule\per\mole\per\kelvin}$
\end{itemize}
Bereken de reactie-enthalpie bij \SI{500}{\celsius} (\SI{773}{\kelvin}). Verklaar waarom de reactie bij hoge temperatuur minder exotherm is dan bij kamertemperatuur.

\item \textbf{Verdampen van water: combinatie van fase-overgang en opwarmen.} \\
Een stoomketel moet \SI{10}{\mole} water per minuut omzetten van vloeibaar water bij \SI{25}{\celsius} naar droge verzadigde stoom bij \SI{100}{\celsius} en \SI{1}{\atm}. Gebruik de volgende gegevens:
\begin{itemize}
\item $C_{p,m}(\text{water, vloeibaar}) \approx \SI{75.3}{\joule\per\mole\per\kelvin}$
\item $\Delta H_{\text{vap}} \approx \SI{40.7}{\kilo\joule\per\mole}$ bij \SI{100}{\celsius}
\end{itemize}
\begin{enumerate}
\item[(a)] Bereken de totale warmte die per minuut moet worden toegevoerd.
\item[(b)] Splits deze op in de bijdrage van opwarmen en de bijdrage van verdampen.
\item[(c)] Wat is de verhouding van beide bijdragen? Wat leert dit over het energie\"erverbruik van een reboiler versus een gewone vloeistofverwarmer?
\end{enumerate}

\item \textbf{Reëel gasgedrag met Van der Waals.} \\
Koolstofdioxide bij hoge druk wijkt sterk af van ideaal gedrag. De Van der Waals-constanten voor $\mathrm{CO_2}$ zijn $a = \SI{0.364}{\pascal\metre\six\per\mole\tothe{2}}$ en $b = \SI{4.27e-5}{\metre\cubed\per\mole}$.
\begin{enumerate}
\item[(a)] Bereken de druk die de vergelijking van Van der Waals voorspelt voor \SI{1}{\mole} $\mathrm{CO_2}$ bij $T = \SI{300}{\kelvin}$ en $V_m = \SI{0.5}{\litre\per\mole}$.
\item[(b)] Vergelijk met de druk die de ideale gaswet voorspelt onder dezelfde omstandigheden.
\item[(c)] Bereken de compressibiliteitsfactor $Z$ in beide gevallen. Verklaar waarom $Z < 1$.
\end{enumerate}

\end{enumerate}

%==============================================================================
% Oefeningen bij Hoofdstuk 2 -- De moleculaire oorsprong van chemische energie
%==============================================================================

\section*{Oefeningen bij Hoofdstuk 2}

\addcontentsline{toc}{section}{Oefeningen bij Hoofdstuk 2}

\begin{enumerate}

\item \textbf{Energie van fotonen en het foto-elektrisch effect.} \\
Cesium heeft een elektronenbindingsenergie (werkfunctie) van \SI{1.95}{\electronvolt}.
\begin{enumerate}
\item[(a)] Bereken de minimale frequentie en golflengte van licht die nodig is om elektronen uit een cesiumoppervlak los te slaan.
\item[(b)] Wordt zichtbaar licht (ca. \SIrange{400}{700}{\nano\metre}) hiervoor gebruikt? Of eerder UV?
\item[(c)] Bereken de maximale kinetische energie van de losgeslagen elektronen wanneer cesium wordt belicht met licht van \(\lambda = \SI{400}{\nano\metre}\).
\end{enumerate}
\emph{Hint:} gebruik $E_{\text{foton}} = h\nu = hc/\lambda$ en $1~\mathrm{eV} = \SI{1.602e-19}{\joule}$.

\item \textbf{Energieniveaus van het waterstofatoom.} \\
De energieniveaus van het waterstofatoom worden gegeven door $E_n = -13{,}6~\mathrm{eV}/n^2$.
\begin{enumerate}
\item[(a)] Bereken de energie van de overgang $n=3 \to n=2$ (Balmer-reeks) en de golflengte van het uitgezonden foton.
\item[(b)] Vergelijk met de overgang $n=2 \to n=1$ (Lyman-reeks). Welke reeks ligt in het zichtbare gebied?
\item[(c)] Hoeveel energie is minimaal nodig om een elektron uit de grondtoestand volledig te ioniseren (tot $n \to \infty$)?
\end{enumerate}

\item \textbf{Bindingsenergie en reactie-enthalpie.} \\
Beschouw de synthese van waterstofchloride uit de elementen:
\[
\mathrm{H_2(g) + Cl_2(g) \rightarrow 2\,HCl(g)}.
\]
Gebruik de volgende gemiddelde bindingsenergie\"en:
\begin{itemize}
\item $D(\mathrm{H{-}H}) = \SI{436}{\kilo\joule\per\mole}$
\item $D(\mathrm{Cl{-}Cl}) = \SI{243}{\kilo\joule\per\mole}$
\item $D(\mathrm{H{-}Cl}) = \SI{431}{\kilo\joule\per\mole}$
\end{itemize}
\begin{enumerate}
\item[(a)] Bereken de reactie-enthalpie uit de bindingsenergie\"en.
\item[(b)] Vergelijk met de waarde die men via standaardvormingsenthalpie\"en zou krijgen ($\Delta H_f^\circ(\mathrm{HCl,g}) = \SI{-92.3}{\kilo\joule\per\mole}$). Verklaar eventuele verschillen.
\item[(c)] Waarom is deze benadering met gemiddelde bindingsenergie\"en slechts een schatting?
\end{enumerate}

\item \textbf{Uitsluitingsprincipe en elektronenconfiguraties.} \\
Geef de volledige elektronenconfiguratie (in termen van orbitalen) van de volgende elementen in hun grondtoestand:
\begin{enumerate}
\item[(a)] koolstof ($Z=6$),
\item[(b)] zuurstof ($Z=8$),
\item[(c)] natrium ($Z=11$).
\end{enumerate}
Verklaar voor zuurstof waarom de twee 2p-elektronen in afzonderlijke orbitalen met parallelle spin worden geplaatst (regel van Hund) en hoe dit verband houdt met het uitsluitingsprincipe van Pauli.

\item \textbf{Activeringsenergie en katalysator.} \\
Een chemische reactie heeft een activeringsenergie van \SI{150}{\kilo\joule\per\mole}. Een katalysator verlaagt deze tot \SI{75}{\kilo\joule\per\mole}.
\begin{enumerate}
\item[(a)] Maak een schematisch energieprofiel (qualitatief) voor de ongekat \-alyseerde en de gekatalyseerde reactie.
\item[(b)] Wat verandert er aan de reactie-enthalpie $\Delta_r H$ door de katalysator? En aan de evenwichtsconstante?
\item[(c)] Hoeveel keer sneller wordt de reactie bij kamertemperatuur (\SI{298}{\kelvin}) dankzij de katalysator? Gebruik de Arrhenius-benadering $k \propto e^{-E_a/RT}$.
\end{enumerate}
\emph{Hint:} bereken de verhouding $k_{\text{kat}}/k_{\text{onkat}} = \exp\!\bigl((E_{a,\text{onkat}} - E_{a,\text{kat}})/RT\bigr)$.

\item \textbf{Constructieve en destructieve interferentie.} \\
Wanneer twee golffuncties $\psi_1$ en $\psi_2$ worden gecombineerd, ontstaat $\psi_{\pm} = \psi_1 \pm \psi_2$.
\begin{enumerate}
\item[(a)] Schets kwalitatief de waarschijnlijkheidsverdeling $|\psi_+|^2$ en $|\psi_-|^2$ in het gebied tussen twee kernen.
\item[(b)] Welke van beide correspondeert met een bindende toestand, welke met een antibindende? Verklaar vanuit de gemiddelde kinetische energie.
\item[(c)] Waarom wordt in de chemie vrijwel uitsluitend met de bindende toestand gewerkt?
\end{enumerate}

\item \textbf{Brug naar thermodynamica: schatting van reactie-enthalpie.} \\
De standaardvormingsenthalpie van waterstofsulfide bedraagt $\Delta H_f^\circ(\mathrm{H_2S,g}) = \SI{-20.6}{\kilo\joule\per\mole}$. De bindingsenergie van een S--H binding bedraagt \SI{363}{\kilo\joule\per\mole} en die van een H--H binding \SI{436}{\kilo\joule\per\mole}. Schat hieruit de bindingsenergie van een S--S binding in $\mathrm{H_2S_2}$, gegeven
\[
\mathrm{2\,H_2S(g) \rightarrow H_2S_2(g) + H_2(g)} \qquad \Delta_r H^\circ \approx \SI{+10}{\kilo\joule\per\mole}.
\]
\emph{Hint:} aan de kant van reagentia worden twee S--H bindingen verbroken; aan de kant van producten worden \'e\'en S--S en \'e\'en H--H binding gevormd.

\end{enumerate}

%==============================================================================
% Oefeningen bij hoofdstuk intermoleculaire
%==============================================================================

\section*{Oefeningen bij Hoofdstuk Intermoleculaire krachten}
\addcontentsline{toc}{section}{Oefeningen}

\begin{enumerate}

\item \textbf{Kookpuntvergelijking van homologe reeksen.} \\
De kookpunten van de $n$-alkanen bij \SI{1}{\atm} zijn:
\begin{itemize}
\item $\mathrm{CH_4}$: \SI{-162}{\celsius}
\item $\mathrm{C_2H_6}$: \SI{-89}{\celsius}
\item $\mathrm{C_3H_8}$: \SI{-42}{\celsius}
\item $\mathrm{C_4H_{10}}$: \SI{0}{\celsius}
\item $\mathrm{C_5H_{12}}$: \SI{36}{\celsius}
\item $\mathrm{C_6H_{14}}$: \SI{69}{\celsius}
\end{itemize}
\begin{enumerate}
\item[(a)] Verklaar de trend vanuit intermoleculaire krachten. Welke kracht is dominant?
\item[(b)] Waarom stijgt het kookpunt minder snel naarmate de ketenlanger wordt?
\item[(c)] Voorspel het kookpunt van $n$-octaan ($\mathrm{C_8H_{18}}$).
\end{enumerate}

\item \textbf{Waterstofbruggen: vergelijking van $\mathrm{H_2O}$, $\mathrm{H_2S}$ en $\mathrm{H_2Se}$.} \\
De kookpunten van de chalcogeniden zijn:
\begin{itemize}
\item $\mathrm{H_2O}$: \SI{100}{\celsius}
\item $\mathrm{H_2S}$: \SI{-60}{\celsius}
\item $\mathrm{H_2Se}$: \SI{-42}{\celsius}
\item $\mathrm{H_2Te}$: \SI{-2}{\celsius}
\end{itemize}
\begin{enumerate}
\item[(a)] Verklaar waarom de kookpunt trend stijgt van $\mathrm{H_2S}$ naar $\mathrm{H_2Te}$, maar $\mathrm{H_2O}$ een anomalie vormt.
\item[(b)] Hoeveel waterstofbruggen kan een watermolecule gemiddeld vormen? Verklaar het nettowerk-karakter.
\item[(c)] Bereken de verhouding $\Delta H_{\text{vap}}(\mathrm{H_2O})/\Delta H_{\text{vap}}(\mathrm{H_2S})$. Verklaar de verhouding vanuit de aard van de intermoleculaire krachten.
\end{enumerate}

\item \textbf{Lennard-Jones: parameters en kookpunt.} \\
Voor methaan gelden $\sigma = \SI{0.381}{\nano\metre}$, $\varepsilon/k_B = \SI{148}{\kelvin}$.
\begin{enumerate}
\item[(a)] Bereken de evenwichtsafstand $r_{\text{eq}}$.
\item[(b)] Bereken de molair bindingsenergie $N_A \varepsilon$.
\item[(c)] Vergelijk met de experimentele verdampingsenthalpie $\Delta H_{\text{vap}} = \SI{8.2}{\kilo\joule\per\mole}$. Verklaar de verhouding (de molecule heeft ongeveer 12 buren in de vloeistof).
\end{enumerate}

\item \textbf{Hildebrand-parameter en mengbaarheid.} \\
Gegeven de Hildebrand-parameters:
\begin{itemize}
\item Tolueen: $\delta = \SI{18.2}{\mega\pascal^{1/2}}$
\item Cyclohexaan: $\delta = \SI{16.8}{\mega\pascal^{1/2}}$
\item Methanol: $\delta = \SI{29.7}{\mega\pascal^{1/2}}$
\item Water: $\delta = \SI{47.9}{\mega\pascal^{1/2}}$
\end{itemize}
\begin{enumerate}
\item[(a)] Welke paar vloeistoffen zullen naar verwachting volledig mengen? Welke niet?
\item[(b)] Bereken de mengenthalpie per liter voor een 50/50-mengsel van tolueen en methanol.
\item[(c)] Waarom wordt de Hildebrand-parameter niet goed toegepast op waterige systemen?
\end{enumerate}

\item \textbf{Brug naar Van der Waals: van $\varepsilon$ naar $a$.} \\
De Van der Waals-parameter $a$ kan worden geschat als $a \approx \tfrac{4}{3} \pi N_A^2 \varepsilon \sigma^3$.
\begin{enumerate}
\item[(a)] Bereken $a$ voor argon ($\sigma = \SI{0.340}{\nano\metre}$, $\varepsilon = \SI{1.65e-21}{\joule}$).
\item[(b)] Vergelijk met de experimentele waarde $a = \SI{0.136}{\pascal\metre\six\per\mole\squared}. Verklaar eventuele verschillen.
\item[(c)] Verklaar waarom de parameter $a$ groter is voor $\mathrm{CO_2}$ dan voor $\mathrm{N_2}$.
\end{enumerate}

\item \textbf{Polariseerbaarheid en dispersie: ordening van halogenen.} \\
De polariseerbaarheden (in \si{\angstrom\cubed}) zijn:
\begin{itemize}
\item $\mathrm{F_2}$: 1{,}56
\item $\mathrm{Cl_2}$: 4{,}61
\item $\mathrm{Br_2}$: 7{,}15
\item $\mathrm{I_2}$: 10{,}5
\end{itemize}
\begin{enumerate}
\item[(a)] Verklaar de trend vanuit de elektronenstructuur.
\item[(b)] Bereken de verhouding van London-dispersie-energie tussen $\mathrm{I_2}$ en $\mathrm{F_2}$ (verondersteld gelijke ionisatienegi\"e en afstand).
\item[(c)] Verbind met de kookpunt-trend van de halogenen.
\end{enumerate}

\item \textbf{Ge\"integreerde case: oplosmiddelselectie voor extractie.} \\
Een farmaceutisch bedrijf moet een actief ingredi\"ent (API) uit een waterige oplossing extraheren. De API heeft een Hildebrand-parameter $\delta_{\text{API}} = \SI{21.0}{\mega\pascal^{1/2}}$. Drie kandidaat-oplosmiddelen worden overwogen:
\begin{itemize}
\item Ethylacetaat: $\delta = \SI{18.6}{\mega\pascal^{1/2}}$
\item Dichloormethaan: $\delta = \SI{20.2}{\mega\pascal^{1/2}}$
\item Hexaan: $\delta = \SI{14.9}{\mega\pascal^{1/2}}$
\end{itemize}
\begin{enumerate}
\item[(a)] Welk oplosmiddel biedt de hoogste oplosbaarheid voor de API?
\item[(b)] Welke oplosmiddel zal zich het best scheiden van de waterige fase? (Verbind met intermoleculaire krachten.)
\item[(c)] Welke bijkomende overwegingen spelen een rol in de praktijk (veiligheid, toxiciteit, milieubelasting)?
\end{enumerate}

\end{enumerate}

%==============================================================================
% Oefeningen bij Hoofdstuk 3 -- Spontaniteit, entropie en Gibbs-vrije energie
%==============================================================================

\section*{Oefeningen bij Hoofdstuk 3}

\addcontentsline{toc}{section}{Oefeningen bij Hoofdstuk 3}

\begin{enumerate}

\item \textbf{Mengentropie van een driecomponent-mengsel.} \\
Een vat bevat aanvankelijk \SI{1}{\mole} stikstof, \SI{1}{\mole} zuurstof en \SI{1}{\mole} argon, elk in afzonderlijke compartimenten bij gelijke temperatuur en druk. De compartimenten worden geopend en de gassen mengen spontaan.
\begin{enumerate}
\item[(a)] Bereken de mengentropie per mol mengsel.
\item[(b)] Bereken de bijbehorende verandering van de Gibbs-vrije energie bij \SI{298}{\kelvin}.
\item[(c)] Hoeveel minimale arbeid is nodig om het mengsel bij \SI{298}{\kelvin} terug te scheiden in de drie zuivere componenten?
\end{enumerate}
\emph{Hint:} $\Delta S_{\text{meng}} = -R\sum_i x_i \ln x_i$; $\Delta G_{\text{meng}} = -T \Delta S_{\text{meng}}$ voor ideale mengsels.

\item \textbf{Entropie bij faseovergangen.} \\
Thallium smelt bij \SI{304}{\celsius} met een smeltenthalpie van $\Delta H_{\text{fus}} = \SI{4.27}{\kilo\joule\per\mole}$.
\begin{enumerate}
\item[(a)] Bereken de smeltentropie.
\item[(b)] Verklaar vanuit de orde-grootte waarom deze entropie kleiner is dan de verdampingsentropie van water ($\Delta S_{\text{vap}} \approx \SI{109}{\joule\per\mole\per\kelvin}$).
\item[(c)] Bereken de verandering van de Gibbs-vrije energie bij het smelten precies op het smeltpunt.
\end{enumerate}

\item \textbf{Spontaniteitscriterium: smelten van ijs bij verschillende temperaturen.} \\
Voor het smelten van ijs geldt $\Delta H_{\text{fus}} = \SI{6.0}{\kilo\joule\per\mole}$ en $\Delta S_{\text{fus}} = \SI{22.0}{\joule\per\mole\per\kelvin}$ (verondersteld constant).
\begin{enumerate}
\item[(a)] Bereken $\Delta G$ bij \SI{-10}{\celsius}, \SI{0}{\celsius} en \SI{+10}{\celsius}.
\item[(b)] Bepaal bij welke temperatuur ijs en vloeibaar water in evenwicht zijn.
\item[(c)] Verklaar vanuit het teken van $\Delta G$ waarom ijs bij \SI{-10}{\celsius} niet spontaan smelt.
\end{enumerate}

\item \textbf{Evenwichtsconstante en standaardreactie-Gibbsenergie.} \\
Voor de reactie
\[
\mathrm{N_2O_4(g) \rightleftharpoons 2\,NO_2(g)}
\]
bedraagt $\Delta_r G^\circ = \SI{+4.8}{\kilo\joule\per\mole}$ bij \SI{298}{\kelvin}.
\begin{enumerate}
\item[(a)] Bereken de evenwichtsconstante $K$ bij \SI{298}{\kelvin}.
\item[(b)] Is de reactie bij standaardomstandigheden spontaan in de voorwaartse richting?
\item[(c)] Bij welke waarde van het reactiequoti\"ent $Q$ zou de reactie precies in evenwicht zijn?
\end{enumerate}

\item \textbf{Invloed van temperatuur: de vergelijking van Van 't Hoff.} \\
Voor de ammoniaksynthese $\mathrm{N_2 + 3\,H_2 \rightleftharpoons 2\,NH_3}$ bedraagt $\Delta_r H^\circ = \SI{-92.4}{\kilo\joule\per\mole}$ en $\Delta_r G^\circ = \SI{-33.0}{\kilo\joule\per\mole}$ bij \SI{298}{\kelvin}.
\begin{enumerate}
\item[(a)] Bereken $K$ bij \SI{298}{\kelvin}.
\item[(b)] Schat $K$ bij \SI{700}{\kelvin} (veronderstel $\Delta_r H^\circ$ constant). Gebruik de ge\"integreerde vorm van Van 't Hoff.
\item[(c)] Verklaar vanuit het teken van $\Delta_r H^\circ$ waarom de evenwichtsconstante daalt bij stijgende temperatuur.
\item[(d)] Combineer met Hoofdstuk 1: bereken ook $\Delta_r S^\circ$ bij \SI{298}{\kelvin}.
\end{enumerate}

\item \textbf{Le Ch\^atelier: invloed van druk op ammoniaksynthese.} \\
Beschouw opnieuw de ammoniaksynthese. Bij een bepaalde temperatuur bedraagt $K = \num{1.0e-3}$.
\begin{enumerate}
\item[(a)] Een reactiemengsel bevat aanvankelijk gelijke molfracties $\mathrm{N_2}$, $\mathrm{H_2}$ en $\mathrm{NH_3}$ bij een totale druk van \SI{1}{\bar}. Bereken $Q$ en bepaal de spontane richting.
\item[(b)] Hoe verandert $Q$ wanneer de totale druk wordt verhoogd tot \SI{200}{\bar}? Verklaar waarom een drukverhoging de vorming van ammoniak bevordert.
\item[(c)] Waarom verandert de evenwichtsconstante $K$ zelf \emph{niet} met de druk?
\end{enumerate}
\emph{Hint:} bij een gasreactie met afname van aantal mol geldt $Q \propto 1/p^{\Delta \nu}$.

\item \textbf{Geïntegreerde case: Haber--Bosch bij twee temperaturen.} \\
Een ingenieur moet een ammoniakreactor ontwerpen en overweegt twee bedrijfs temperaturen: \SI{400}{\celsius} en \SI{500}{\celsius}. Voor de synthese $\mathrm{N_2 + 3\,H_2 \rightleftharpoons 2\,NH_3}$ geldt bij \SI{298}{\kelvin}: $\Delta_r H^\circ = \SI{-92.4}{\kilo\joule\per\mole}$, $\Delta_r G^\circ = \SI{-33.0}{\kilo\joule\per\mole}$, en $\Delta_r C_p^\circ \approx \SI{-5.2}{\joule\per\mole\per\kelvin}$ (constant verondersteld).
\begin{enumerate}
\item[(a)] Gebruik de wet van Kirchhoff (Hoofdstuk 1) om $\Delta_r H^\circ$ te berekenen bij beide temperaturen.
\item[(b)] Bereken $\Delta_r S^\circ$ bij \SI{298}{\kelvin} en schat $\Delta_r S^\circ$ bij \SI{500}{\celsius} via $\Delta_r S^\circ(T) = \Delta_r S^\circ(T_0) + \Delta_r C_p^\circ \ln(T/T_0)$.
\item[(c)] Bereken $\Delta_r G^\circ$ bij beide temperaturen via $G = H - TS$.
\item[(d)] Bereken $K$ bij beide temperaturen.
\item[(e)] Concludeer: welke temperatuur levert thermodynamisch de hoogste opbrengst? Leg uit waarom de industrie toch bij \SIrange{400}{500}{\celsius} werkt.
\end{enumerate}
\emph{Hint:} deze oefening combineert alle hoofdstukken 1--3: Kirchhoff (H1), enthalpie/entropie/Gibbs (H3), en de brug met de evenwichtsconstante (H3).

\end{enumerate}


%==============================================================================
% Oefeningen bij Hoofdstuk 4
%==============================================================================

\section*{Oefeningen bij Hoofdstuk 4}
\addcontentsline{toc}{section}{Oefeningen bij Hoofdstuk 4}

\begin{enumerate}

\item \textbf{Dampdruk via Antoine.} \\
Gebruik de Antoine-vergelijking om de dampdruk van water te berekenen bij \SI{50}{\celsius}, \SI{80}{\celsius} en \SI{120}{\celsius}. Neem de constanten $A = 8{,}07131$, $B = 1730{,}63$, $C = 233{,}426$ ($p$ in \si{\milli\metre\of\mercury}, $T$ in \si{\celsius}, geldig \SIrange{1}{100}{\celsius}). Vergelijk de waarde bij \SI{50}{\celsius} met de schatting uit Clausius--Clapeyron (Oefening 3 van H1). Verklaar eventuele verschillen.

\item \textbf{Bubbelpunt en dampsamenstelling (Raoult).} \\
Een binair mengsel van $n$-hexaan (A) en $n$-octaan (B) met $x_A = 0{,}30$ bevindt zich bij \SI{80}{\celsius}. De verzadigingsdampdrukken bedragen $p_A^* = \SI{1.84}{\bar}$ en $p_B^* = \SI{0.32}{\bar}$.
\begin{enumerate}
\item[(a)] Bereken de bubbelpuntdruk.
\item[(b)] Bereken de dampsamenstelling $y_A$, $y_B$.
\item[(c)] Bereken de relatieve vluchtigheid $\alpha_{AB}$ en beoordeel de haalbaarheid van destillatie.
\end{enumerate}

\item \textbf{Hefboomwet.} \\
Voor het mengsel uit Oefening 2, bij de bubbelpuntdruk, wordt de globale samenstelling $z_A = 0{,}50$ gekozen. Bereken de verhouding $n_L/n_V$ en de absolute molhoeveelheden vloeistof en damp per mol mengsel.

\item \textbf{Azeotroop-herkenning.} \\
Beschrijf of de volgende binaire mengsels een minimaal of maximaal kokende azeotroop vormen, en leg uit vanuit de intermoleculaire interacties:
\begin{enumerate}
\item[(a)] water--ethanol,
\item[(b)] aceton--chloroform,
\item[(c)] salpeterzuur--water.
\end{enumerate}

\item \textbf{Activiteitsco\"effici\"ent uit dampdrukmeting.} \\
Voor een binair mengsel van aceton (A) en chloroform (B) bij \SI{35}{\celsius} wordt bij $x_A = 0{,}3$ een parti\"ele dampdruk $p_A = \SI{0.20}{\bar}$ gemeten. De verzadigingsdampdruk van zuiver aceton bedraagt $p_A^* = \SI{0.46}{\bar}$.
\begin{enumerate}
\item[(a)] Bereken de activiteitsco\"effici\"ent $\gamma_A$.
\item[(b)] Is dit een positieve of negatieve afwijking? Verklaar.
\item[(c)] Schat de parameter $A$ in een Margules-\'e\'enparameter-model ($\ln \gamma_A = A x_B^2$).
\end{enumerate}

\item \textbf{Vloeistof-vloeistofextractie met Nernst.} \\
Een waterige oplossing bevat \SI{0.05}{\mole\per\litre} benzo\"ezzuur. De verdelingsco\"effici\"ent tussen tolueen en water bedraagt $K_D = 12$ (ten gunste van tolueen). \SI{1}{\litre} waterige oplossing wordt ge\"extraheerd met \SI{0.5}{\litre} tolueen.
\begin{enumerate}
\item[(a)] Bereken de resterende concentratie in de waterige fase na extractie.
\item[(b)] Hoeveel extractiestappen met vers tolueen zijn nodig om \SI{99}{\percent} te verwijderen?
\end{enumerate}

\item \textbf{Osmotische druk van zeewater.} \\
Zeewater bevat ongeveer \SI{0.55}{\mole\per\litre} opgeloste zouten (voornamelijk NaCl, volledig gedissocieerd).
\begin{enumerate}
\item[(a)] Bereken de osmotische druk bij \SI{25}{\celsius}.
\item[(b)] Welke minimale druk moet een omgekeerde-osmosemembraan leveren om zuiver water te produceren?
\item[(c)] Hoeveel arbeid is minimaal nodig om \SI{1}{\litre} zuiver water uit zeewater te produceren?
\end{enumerate}

\end{enumerate}

%==============================================================================
% Oefeningen bij Hoofdstuk 5
%==============================================================================

\section*{Oefeningen bij Hoofdstuk 5}
\addcontentsline{toc}{section}{Oefeningen bij Hoofdstuk 5}

\begin{enumerate}

\item \textbf{Evenwichtsconstante uit $\Delta_r G^\circ$.} \\
Voor de reactie $\mathrm{N_2O_4(g) \rightleftharpoons 2\,NO_2(g)}$ geldt bij \SI{298}{\kelvin}: $\Delta_r G^\circ = \SI{+4.8}{\kilo\joule\per\mole}$.
\begin{enumerate}
\item[(a)] Bereken $K_p$ bij \SI{298}{\kelvin}.
\item[(b)] Is de reactie bij standaardomstandigheden spontaan in de voorwaartse richting?
\item[(c)] Bij welke waarde van $Q$ geldt $\Delta_r G = 0$?
\end{enumerate}

\item \textbf{Verband $K_p$ en $K_c$.} \\
Voor $\mathrm{N_2 + 3\,H_2 \rightleftharpoons 2\,NH_3}$ geldt $K_p = 1{,}5 \times 10^{-4}$ bij \SI{700}{\kelvin}. Bereken $K_c$ bij dezelfde temperatuur. Verklaar waarom $K_c \neq K_p$.

\item \textbf{Reactietabel (ICE-table) voor ammoniak.} \\
Beschouw de ammoniaksynthese bij $T = \SI{700}{\kelvin}$, $P = \SI{100}{\bar}$, $K_p = 1{,}5 \times 10^{-4}$. De beginhoeveelheden bedragen $n_{\mathrm{N_2}}^0 = 1$, $n_{\mathrm{H_2}}^0 = 3$, $n_{\mathrm{NH_3}}^0 = 0$ mol.
\begin{enumerate}
\item[(a)] Stel de reactietabel op.
\item[(b)] Stel de evenwichtsvergelijking in $\xi$ op.
\item[(c)] Los numeriek op en bereken de omzetting van stikstof.
\item[(d)] Hoe verandert de omzetting wanneer \SI{2}{\mole} argon (inert) wordt toegevoegd bij constante totale druk?
\end{enumerate}

\item \textbf{Heterogeen evenwicht: ontleding van kalksteen.} \\
Voor $\mathrm{CaCO_3(s) \rightleftharpoons CaO(s) + CO_2(g)}$ geldt bij \SI{800}{\celsius} $K_p = 0{,}22$.
\begin{enumerate}
\item[(a)] Wat is de evenwichtsdruk van $\mathrm{CO_2}$?
\item[(b)] Bij welke omgevingsdruk van $\mathrm{CO_2}$ zal kalksteen spontaan ontleden?
\item[(c)] Waarom verschijnen de vaste fasen niet in de evenwichtsconstante?
\end{enumerate}

\item \textbf{Oplosbaarheidsproduct.} \\
Voor $\mathrm{Ag_2CrO_4(s) \rightleftharpoons 2\,Ag^+(aq) + CrO_4^{2-}(aq)}$ geldt $K_{ps} = 1{,}1 \times 10^{-12}$ bij \SI{25}{\celsius}.
\begin{enumerate}
\item[(a)] Bereken de oplosbaarheid $s$ in zuiver water.
\item[(b)] Hoe verandert $s$ wanneer \SI{0.1}{\mole\per\litre} $\mathrm{AgNO_3}$ wordt toegevoegd (gemeenschappelijk-ioneneffect)?
\end{enumerate}

\item \textbf{Gekoppelde evenwichten: syngas.} \\
Beschouw stoomreforming van methaan met de twee simultane reacties:
\[
\mathrm{CH_4 + H_2O \rightleftharpoons CO + 3\,H_2} \quad (K_1),
\]
\[
\mathrm{CO + H_2O \rightleftharpoons CO_2 + H_2} \quad (K_2).
\]
\begin{enumerate}
\item[(a)] Welke componenten zijn gemeenschappelijk?
\item[(b)] Hoeveel onafhankelijke evenwichtsvergelijkingen moeten worden opgelost?
\item[(c)] Waarom kan men niet zomaar de twee reacties afzonderlijk behandelen?
\end{enumerate}

\item \textbf{Ge\"integreerde case: methanolproductie.} \\
Methanol wordt industrieel geproduceerd via:
\[
\mathrm{CO(g) + 2\,H_2(g) \rightleftharpoons CH_3OH(g)},
\]
met $\Delta_r H^\circ = \SI{-90.7}{\kilo\joule\per\mole}$ en $\Delta_r G^\circ = \SI{-25.3}{\kilo\joule\per\mole}$ bij \SI{298}{\kelvin}.
\begin{enumerate}
\item[(a)] Bereken $K_p$ bij \SI{298}{\kelvin}.
\item[(b)] Is de reactie exotherm of endotherm? Wat impliceert dit voor de temperatuurkeuze?
\item[(c)] Hoe verandert de molfractie-matige evenwichtsconstante $K_x$ met de totale druk? Verklaar waarom industri\"ele methanolproductie bij hoge druk (\SIrange{50}{100}{\bar}) werkt.
\item[(d)] Schat $K_p$ bij \SI{500}{\kelvin} via Van 't Hoff ($\Delta_r H^\circ$ constant verondersteld).
\item[(e)] Waarom wordt in de praktijk bij \SIrange{200}{300}{\celsius} gewerkt, ondanks de thermodynamisch gunstige lagere temperatuur? (Combineer met Hoofdstuk~2: activeringsenergie en katalysator.)
\end{enumerate}

\end{enumerate}

%==============================================================================
% Oefeningen bij Hoofdstuk 6
%==============================================================================

\section*{Oefeningen bij Hoofdstuk 6}
\addcontentsline{toc}{section}{Oefeningen bij Hoofdstuk 6}

\begin{enumerate}

\item \textbf{Celpotentiaal en $\Delta_r G^\circ$.} \\
Voor de reactie $\mathrm{Zn + Cu^{2+} \rightarrow Zn^{2+} + Cu}$ geldt $E^\circ = \SI{1.10}{\volt}$.
\begin{enumerate}
\item[(a)] Bereken $\Delta_r G^\circ$.
\item[(b)] Bereken de evenwichtsconstante $K$ bij \SI{298}{\kelvin}.
\item[(c)] Is de reactie spontaan bij standaardomstandigheden?
\end{enumerate}

\item \textbf{Nernstvergelijking voor de Daniell-cel.} \\
Beschouw de Daniell-cel met $[\mathrm{Zn^{2+}}] = \SI{0.10}{\mole\per\litre}$ en $[\mathrm{Cu^{2+}}] = \SI{1.0}{\mole\per\litre}$.
\begin{enumerate}
\item[(a)] Bereken de celpotentiaal bij \SI{298}{\kelvin}.
\item[(b)] Hoe verandert $E$ wanneer beide concentraties worden gehalveerd?
\end{enumerate}

\item \textbf{Elektrolyse van water --- minimumspanning.} \\
Voor $\mathrm{2\,H_2O \rightarrow 2\,H_2 + O_2}$ geldt $\Delta_r G^\circ = \SI{+474.4}{\kilo\joule\per\mole}$ bij \SI{298}{\kelvin}.
\begin{enumerate}
\item[(a)] Bereken de thermodynamische minimumspanning.
\item[(b)] Bereken de elektrische arbeid per mol $\mathrm{H_2}$.
\item[(c)] Een industri\"ele elektrolyseur werkt bij \SI{1.9}{\volt}. Bereken de voltageffici\"entie.
\end{enumerate}

\item \textbf{Brandstofcel: maximale arbeid.} \\
Een PEM-brandstofcel oxydeert \SI{1}{\mole} $\mathrm{H_2}$ met $\mathrm{O_2}$ tot vloeibaar water bij \SI{298}{\kelvin}.
\begin{enumerate}
\item[(a)] Bereken de maximale elektrische arbeid per mol $\mathrm{H_2}$.
\item[(b)] Bereken de standaardcelpotentiaal.
\item[(c)] Een cel levert \SI{0.7}{\volt} bij \SI{0.5}{\ampere}. Bereken het thermodynamisch rendement.
\end{enumerate}

\item \textbf{Pourbaix-diagram van ijzer interpreteren.} \\
Beschouw het Pourbaix-diagram van ijzer.
\begin{enumerate}
\item[(a)] Teken schematisch de drie gebieden (immuniteit, corrosie, passiviteit) in een potentiaal--pH-vlak.
\item[(b)] Waarom bevindt ijzer bij lage pH en hoge potentiaal zich in het corrosiegebied?
\item[(c)] Hoe kan kathodische bescherming de thermodynamische stabiliteit be\"invloeden?
\end{enumerate}

\item \textbf{Materiaalkeuze: aluminium versus ijzer.} \\
Aluminium bezit $E^\circ = \SI{-1.66}{\volt}$ en is dus thermodynamisch veel reactiever dan ijzer ($E^\circ = \SI{-0.44}{\volt}$). Toch wordt aluminium in de luchtvaart veel gebruikt omwille van zijn corrosiebestendigheid.
\begin{enumerate}
\item[(a)] Verklaar dit schijnbare paradox vanuit passiviteit.
\item[(b)] Waarom faalt deze bescherming in sterke zuren of basen?
\item[(c)] Hoe kan men ijzer ``leren'' van aluminium door legering?
\end{enumerate}

\item \textbf{Ge\"integreerde case: groene waterstofproductie.} \\
Een elektrolyseur produceert waterstof door elektrolyse van water bij \SI{80}{\celsius} en \SI{30}{\bar}. De elektrische energie\"erinvoer bedraagt \SI{52.5}{\kilo\watthour} per kg $\mathrm{H_2}$.
\begin{enumerate}
\item[(a)] Bereken de thermodynamische minimumarbeid per kg $\mathrm{H_2}$ (gebruik $\Delta_r G^\circ \approx \SI{237}{\kilo\joule\per\mole}$ bij \SI{298}{\kelvin}; neem een benadering via Kirchhoff met $\Delta_r C_p^\circ \approx -\SI{80}{\joule\per\mole\per\kelvin}$ voor \SI{80}{\celsius}).
\item[(b)] Bereken de energie\"ereffici\"entie van de elektrolyseur.
\item[(c)] Hoeveel kg $\mathrm{CO_2}$ wordt vermeden per kg $\mathrm{H_2}$ wanneer de stroom uit hernieuwbare bronnen komt en waterstof een verbrandingsenergie van \SI{142}{\mega\joule\per\kilogram} heeft, vergeleken met verbranding van aardgas met $\Delta H_{\text{verbr}} \approx \SI{55.5}{\mega\joule\per\kilogram}$?
\item[(d)] Combineer met Hoofdstuk~5: hoe zou een chemisch evenwichtanalyse helpen om de electrolyse bij lagere celspanning te ontwerpen?
\end{enumerate}

\end{enumerate}

%==============================================================================
% Oefeningen bij Hoofdstuk 7
%==============================================================================

\section*{Oefeningen bij Hoofdstuk 7}
\addcontentsline{toc}{section}{Oefeningen bij Hoofdstuk 7}

\begin{enumerate}

\item \textbf{Eerste-orde reactie en halfwaardetijd.} \\
Een bepaalde reactie $A \rightarrow P$ heeft $k = \SI{2.5e-3}{\per\second}$ bij \SI{298}{\kelvin}.
\begin{enumerate}
\item[(a)] Bereken de halfwaardetijd.
\item[(b)] Hoe lang duurt het om \SI{99}{\percent} van $A$ om te zetten?
\item[(c)] Bereken de beginconcentratie-indicator $\ln[A]$ na \SI{500}{\second}.
\end{enumerate}

\item \textbf{Tweede-orde reactie.} \\
Voor $2A \rightarrow P$ geldt $k = \SI{0.04}{\litre\per\mole\per\second}$ bij \SI{350}{\kelvin}. De beginconcentratie bedraagt $[A]_0 = \SI{0.50}{\mole\per\litre}$.
\begin{enumerate}
\item[(a)] Bereken de concentratie na \SI{60}{\second}.
\item[(b)] Bereken de halfwaardetijd.
\item[(c)] Hoe lang duurt het om \SI{99}{\percent} om te zetten? Vergelijk met Oefening 1.
\end{enumerate}

\item \textbf{Arrhenius: $E_a$ uit metingen.} \\
Voor een bepaalde reactie zijn de snelheidsconstanten:
\begin{itemize}
\item $k = \SI{1.2e-4}{\per\second}$ bij \SI{300}{\kelvin}
\item $k = \SI{8.4e-3}{\per\second}$ bij \SI{340}{\kelvin}
\end{itemize}
\begin{enumerate}
\item[(a)] Bereken $E_a$.
\item[(b)] Bereken de pre\"exponenti\"ele factor $A$.
\item[(c)] Schat $k$ bij \SI{320}{\kelvin}.
\end{enumerate}

\item \textbf{Steady-state-benadering.} \\
Beschouw het mechanisme:
\[
A \xrightarrow{k_1} B \xrightarrow{k_2} C,
\]
met $k_1 = \SI{0.10}{\per\second}$ en $k_2 = \SI{1.0}{\per\second}$. De beginconcentratie bedraagt $[A]_0 = \SI{1.0}{\mole\per\litre}$, $[B]_0 = [C]_0 = 0$.
\begin{enumerate}
\item[(a)] Pas de steady-state toe op $[B]$ en leid $[B](t)$ af.
\item[(b)] Bepaal $[A](t)$ en $[C](t)$.
\item[(c)] Wanneer is de steady-state-benadering gewettigd?
\end{enumerate}

\item \textbf{Katalysator en Arrhenius.} \\
Een ongekat \-al yseerde reactie heeft $E_{a,\text{onkat}} = \SI{140}{\kilo\joule\per\mole}$, $A = \SI{1.0e13}{\per\second}$. Met een katalysator daalt $E_a$ tot \SI{80}{\kilo\joule\per\mole}; $A$ verandert niet.
\begin{enumerate}
\item[(a)] Bereken de snelheidsverhoging bij \SI{298}{\kelvin} en bij \SI{500}{\kelvin}.
\item[(b)] Bij welke temperatuur zou de ongekat \-al yseerde reactie dezelfde snelheid hebben als de gekatalyseerde bij \SI{298}{\kelvin}?
\item[(c)] Verklaar waarom een katalysator een lagere procestemperatuur mogelijk maakt.
\end{enumerate}

\item \textbf{Langmuir--Hinshelwood en oppervlakte-reactie.} \\
Een heterogeen gekatalyseerde bimoleculaire reactie $A + B \rightarrow P$ op een oppervlak volgt:
\[
r = \frac{k K_A K_B p_A p_B}{(1 + K_A p_A + K_B p_B)^2}.
\]
Met $K_A = K_B = \SI{1.0}{\per\bar}$, $k = \SI{2.0e-5}{\mole\per\second\per\meter\squared}$, $p_A = p_B = p$.
\begin{enumerate}
\item[(a)] Bereken $r$ bij $p = \SI{0.1}{\bar}$ en $p = \SI{10}{\bar}$.
\item[(b)] Bepaal bij welke druk $r$ maximaal is.
\item[(c)] Verklaar kwalitatief waarom $r$ daalt bij hoge druk.
\end{enumerate}

\item \textbf{Michaelis--Menten en enzymverzadiging.} \\
Een enzym heeft $r_{\max} = \SI{2.0e-5}{\mole\per\litre\per\second}$ en $K_M = \SI{5.0e-4}{\mole\per\litre}$.
\begin{enumerate}
\item[(a)] Bereken $r$ bij $[S] = K_M$, $[S] = 10 K_M$ en $[S] = 0{,}1 K_M$.
\item[(b)] Bepaal de reactieorde in elk regime (laag, gemiddeld, hoog $[S]$).
\item[(c)] Bij welke $[S]$ is $r = 0{,}90\,r_{\max}$?
\end{enumerate}

\item \textbf{Reactorontwerp: PFR versus CSTR.} \\
Een eerste-orde reactie met $k = \SI{0.08}{\per\second}$ moet worden uitgevoerd tot een omzetting $X = 0{,}95$.
\begin{enumerate}
\item[(a)] Bereken de vereiste verblijftijd voor een PFR.
\item[(b)] Bereken de vereiste verblijftijd voor een CSTR.
\item[(c)] Bereken het volumebesparing van PFR t.o.v. CSTR voor een debiet van \SI{0.5}{\metre\cubed\per\second}.
\end{enumerate}

\item \textbf{Geïntegreerde case: methanolsynthese.} \\
Methanolsynthese $\mathrm{CO + 2\,H_2 \rightleftharpoons CH_3OH}$ is exotherm ($\Delta_r H^\circ = \SI{-90.7}{\kilo\joule\per\mole}$) en verloopt met $\mathrm{Cu/ZnO/Al_2O_3}$-katalysator bij \SIrange{200}{300}{\celsius} en \SIrange{50}{100}{\bar}.
\begin{enumerate}
\item[(a)] Verklaar vanuit thermodynamica (Hoofdstuk~3 en 5) waarom lage temperatuur en hoge druk gunstig zijn.
\item[(b)] Verklaar vanuit kinetiek waarom men niet beneden \SI{200}{\celsius} werkt, ondanks de thermodynamische gunst.
\item[(c)] Schat de snelheidsverhoging van de $\mathrm{Cu}$-katalysator als zij $E_a$ verlaagt van \SI{180}{\kilo\joule\per\mole} (thermische reactie) naar \SI{80}{\kilo\joule\per\mole} bij \SI{250}{\celsius}.
\item[(d)] Waarom is de katalysatortemperatuur in methanolsynthese lager dan in ammoniaksynthese, ondanks een vergelijkbare thermodynamica? Verbind met Hoofdstuk~2 (activeringsenergie en katalysatorwerking).
\end{enumerate}

\item \textbf{Brug naar thermodynamica: $K = k_1/k_{-1}$.} \\
Voor $\mathrm{A \rightleftharpoons B}$ geldt bij \SI{500}{\kelvin}: $k_1 = \SI{1.5e-2}{\per\second}$, $k_{-1} = \SI{3.0e-3}{\per\second}$.
\begin{enumerate}
\item[(a)] Bereken de evenwichtsconstante $K$.
\item[(b)] Bereken $\Delta_r G^\circ$ bij \SI{500}{\kelvin}.
\item[(c)] Als $E_{a,\text{voorw}} = \SI{80}{\kilo\joule\per\mole}$, bereken dan $E_{a,\text{terug}}$ en vergelijk met $\Delta_r H^\circ$.
\end{enumerate}

\end{enumerate}

%==============================================================================
% Oefeningen bij Hoofdstuk 8
%==============================================================================

\section*{Oefeningen bij Hoofdstuk 8}
\addcontentsline{toc}{section}{Oefeningen bij Hoofdstuk 8}

\begin{enumerate}

\item \textbf{Oppervlaktespanning en capillariteit.} \\
Water ($\sigma = \SI{72}{\milli\newton\per\metre}$, $\rho = \SI{1000}{\kilogram\per\metre\cubed}$) stijgt in een glazen capillair met binnendiameter \SI{0.20}{\milli\metre}. De contacthoek bedraagt $\theta = 0^\circ$ (volledige bevochtiging).
\begin{enumerate}
\item[(a)] Bereken de stijghoogte.
\item[(b)] Wat gebeurt er met de stijghoogte in een capillair met hetzelfde materiaal maar met kwik ($\sigma = \SI{485}{\milli\newton\per\metre}$, $\rho = \SI{13600}{\kilogram\per\metre\cubed}$, $\theta = 140^\circ$)? Verklaar het verschil.
\end{enumerate}

\item \textbf{Young--Laplace in longblaasjes.} \\
Een longblaasje heeft straal $r = \SI{100}{\micro\metre}$. Veronderstel water zonder surfactant ($\sigma = \SI{72}{\milli\newton\per\metre}$).
\begin{enumerate}
\item[(a)] Bereken de binnendruk.
\item[(b)] Bereken de nieuwe binnendruk wanneer de straal afneemt tot \SI{50}{\micro\metre}.
\item[(c)] Verklaar fysiologisch waarom longsurfactant noodzakelijk is.
\end{enumerate}

\item \textbf{Langmuir-adsorptie en reactiekinetiek.} \\
Voor de adsorptie van een substraat $S$ op een enzym-achtig oppervlak geldt $K = \SI{2.0e3}{\per\mole\per\litre}$.
\begin{enumerate}
\item[(a)] Bereken de bezettingsgraad $\theta$ bij $[S] = \SI{0.001}{\mole\per\litre}$ en bij $[S] = \SI{0.01}{\mole\per\litre}$.
\item[(b)] Bij welke concentratie geldt $\theta = 0{,}5$?
\item[(c)] Vergelijk met Michaelis--Menten (Hoofdstuk~7): wat is de overeenkomst tussen $K$ en $K_M$?
\end{enumerate}

\item \textbf{BET en specifieke oppervlakte.} \\
Een katalysatordrager heeft een BET-monolaag-volume van \SI{25}{\centi\metre\cubed\per\gram} (STP) voor $\mathrm{N_2}$-adsorptie.
\begin{enumerate}
\item[(a)] Bereken de specifieke oppervlakte ($\mathrm{N_2}$-molecule-oppervlakte $= \SI{0.162}{\nano\metre\squared}$).
\item[(b)] Hoeveel mol Pt kan op deze drager worden gedispergeerd als 1 atoom Pt per \SI{0.09}{\nano\metre\squared} wordt aangenomen?
\end{enumerate}

\item \textbf{CMC en surfactant-thermodynamica.} \\
Voor een niet-ionische surfactant in water wordt de CMC gemeten als functie van de temperatuur:
\begin{itemize}
\item \SI{20}{\celsius}: $\mathrm{CMC} = \SI{5.0e-5}{\mole\per\litre}$
\item \SI{40}{\celsius}: $\mathrm{CMC} = \SI{7.5e-5}{\mole\per\litre}$
\end{itemize}
\begin{enumerate}
\item[(a)] Verklaar waarom de CMC stijgt met de temperatuur.
\item[(b)] Hoeveel vrije surfactant-moleculen bevinden zich in oplossing bij $c = \SI{1e-3}{\mole\per\litre}$ ver boven de CMC? Verklaar vanuit het concept ``vrije concentratie constant boven CMC''.
\end{enumerate}

\item \textbf{Zeta-potentiaal en coagulatie.} \\
Een collo\"idale klei-dispersie in water heeft een initi\"ele zeta-potentiaal van $\zeta = \SI{-45}{\milli\volt}$. Na toevoeging van \SI{10}{\milli\mole\per\litre} $\mathrm{CaCl_2}$ daalt $|\zeta|$ tot \SI{12}{\milli\volt}.
\begin{enumerate}
\item[(a)] Verklaar waarom $\mathrm{Ca^{2+}}$ de zeta-potentiaal verlaagt (vergelijk met $\mathrm{Na^+}$).
\item[(b)] Voorspel de stabiliteit voor en na toevoeging.
\item[(c)] Waarom wordt $\mathrm{CaCl_2}$ vaak vervangen door $\mathrm{Al_2(SO_4)_3}$ in drinkwaterzuivering?
\end{enumerate}

\item \textbf{Emulsietype en HLB.} \\
Een voedingsfabrikant wil een vloeibare margarine (W/O-emulsie) en een dressingsaus (O/W-emulsie) ontwikkelen.
\begin{enumerate}
\item[(a)] Welke HLB-reeks moet de emulgator hebben voor elk product?
\item[(b)] Verklaar vanuit de structuur van lecithine (HLB $\sim 7$) waarom dit \emph{beide} types emulsie kan stabiliseren, en waarom het \emph{geen} ideale keuze is voor een van beide.
\end{enumerate}

\end{enumerate}

%==============================================================================
% Oefeningen bij Hoofdstuk 9
%==============================================================================

\section*{Oefeningen bij Hoofdstuk 9}
\addcontentsline{toc}{section}{Oefeningen bij Hoofdstuk 9}

\begin{enumerate}

\item \textbf{Fourier: warmteverlies doorheen een muur.} \\
Een stalen muur ($\lambda = \SI{50}{\watt\per\metre\per\kelvin}$, dikte \SI{2}{\centi\metre}) scheidt een reactor at \SI{150}{\celsius} van de buitenlucht at \SI{20}{\celsius}.
\begin{enumerate}
\item[(a)] Bereken de warmtestroom per \si{\metre\squared}.
\item[(b)] De reactor heeft een wandoppervlak van \SI{5}{\metre\squared}. Bereken het totale warmteverlies in \si{\kilo\watt}.
\item[(c)] Hoeveel extra isolatie ($\lambda_{\text{isol}} = \SI{0.04}{\watt\per\metre\per\kelvin}$) is nodig om het verlies te halveren?
\end{enumerate}

\item \textbf{Diffusie: tijdschaal in een katalysatorpore.} \\
Een katalysatorpore heeft diameter \SI{10}{\nano\metre}. Reactant A diffundeert in de pore met $D_{\text{eff}} = \SI{1e-8}{\metre\squared\per\second}$ (effectieve diffusieco\"effici\"ent in poreuze medium).
\begin{enumerate}
\item[(a)] Bereken de karakteristieke diffusietijd over \SI{1}{\micro\metre} (typische pore-lengte).
\item[(b)] Hoe verandert deze tijd als de pore-lengte verdubbelt?
\item[(c)] Waarom is interne diffusie vaak de snelheidsbepalende stap in poreuze katalysatoren? (Verbind met Hoofdstuk~7 en H8.)
\end{enumerate}

\item \textbf{Stokes--Einstein: diffusie van eiwitten.} \\
Bereken de diffusieco\"effici\"ent van hemoglobine ($r = \SI{3.0}{\nano\metre}$) in water ($\eta = \SI{1.0e-3}{\pascal\second}$) bij \SI{310}{\kelvin}.
\begin{enumerate}
\item[(a)] Bereken $D$.
\item[(b)] Hoe lang duurt diffusie over \SI{10}{\micro\metre}?
\item[(c)] Vergelijk met de waarde van albumine uit het voorbeeld in het hoofdstuk. Verklaar de verschillen.
\end{enumerate}

\item \textbf{Diffusie-gelimiteerde reactie.} \\
De reactie $\mathrm{A + B \rightarrow P}$ in water heeft $D_A = D_B = \SI{5e-10}{\metre\squared\per\second}$, $r_A + r_B = \SI{0.6}{\nano\metre}$.
\begin{enumerate}
\item[(a)] Bereken de diffusie-gelimiteerde snelheidsconstante $k_D$ via Smoluchowski.
\item[(b)] Een enzym heeft $k_{\text{cat}}/K_M = \SI{5e7}{\litre\per\mole\per\second}. Is dit enzym ``diffusie-perfect''?
\end{enumerate}

\item \textbf{Viscositeit en moleculaire structuur.} \\
\begin{enumerate}
\item[(a)] Verklaar waarom glycerol ($\eta \approx \SI{1.0}{\pascal\second}$) veel visqueuzer is dan water ($\eta \approx \SI{1e-3}{\pascal\second}$) bij kamertemperatuur. (Verbind met waterstofbruggen.)
\item[(b)] Verklaar waarom viscositeit van gassen toeneemt met temperatuur, terwijl die van vloeistoffen afneemt.
\end{enumerate}

\item \textbf{Grenslaag en massaoverdracht.} \\
Een gas $A$ wordt geabsorbeerd in een vloeistof in een packed column. De grenslaagdikte bedraagt $\delta = \SI{50}{\micro\metre}$, $D = \SI{2e-9}{\metre\squared\per\second}$.
\begin{enumerate}
\item[(a)] Bereken de massaoverdrachtsco\"effici\"ent $k_L$.
\item[(b)] Bereken de flux wanneer $c_{\text{gas}} = \SI{0.05}{\mole\per\litre}$ en $c_{\text{vloeistof}} = \SI{0.01}{\mole\per\litre}$.
\end{enumerate}

\item \textbf{Ge\"integreerde case: oxygen transport in fermentoren.} \\
Een bioreactor van \SI{100}{\metre\cubed} bevat bacteri\"ele biomass die \SI{2.0}{\mole\per\second} $\mathrm{O_2}$ verbruikt. Lucht wordt via een sparger ge\"injecteerd; de massaoverdrachtsco\"effici\"ent $k_L a$ bedraagt \SI{0.1}{\per\second} ($a$ is specifiek grensvlakoppervlak).
\begin{enumerate}
\item[(a)] Bereken de oplosbaarheid van $\mathrm{O_2}$ in water bij \SI{1}{\bar} lucht: $c^* = \SI{0.27}{\milli\mole\per\litre}$ (gebruik wet van Henry, H $\sim \SI{4.0e4}{\bar}$).
\item[(b)] Bereken de overdrachtsnelheid per reactorvolume als de bulkconcentratie \SI{0.02}{\milli\mole\per\litre} bedraagt.
\item[(c)] Is de oxygenoverdracht beperkend voor de fermentatie? Verbind met H7 (kinetiek) en H8 (grensvlakken).
\end{enumerate}

\end{enumerate}

%==============================================================================
% Oefeningen bij Hoofdstuk 10
%==============================================================================

\section*{Oefeningen bij Hoofdstuk 10}
\addcontentsline{toc}{section}{Oefeningen bij Hoofdstuk 10}

\begin{enumerate}

\item \textbf{Atomen per eenheidscel en pakking.} \\
Bereken het aantal atomen per eenheidscel en de pakkingsfractie voor:
\begin{enumerate}
\item[(a)] BCC (bv. Cr, $a = \SI{0.291}{\nano\metre}$).
\item[(b)] FCC (bv. Al, $a = \SI{0.405}{\nano\metre}$).
\item[(c)] Primitief kubisch (Polonium, $a = \SI{0.336}{\nano\metre}$).
\end{enumerate}
Verklaar waarom de pakkingsfractie het hoogst is voor FCC.

\item \textbf{Brag: identificatie van een onbekende stof.} \\
Een poederdiffractiepatroon, gemeten met Cu K$\alpha$ ($\lambda = \SI{0.154}{\nano\metre}$), toont een eerste piek bij $\theta = 28{,}3^\circ$.
\begin{enumerate}
\item[(a)] Bereken de intervlakafstand $d$.
\item[(b)] Als de stof Si is met diamantstructuur ($d_{111} = a/\sqrt{3}$), bereken $a$.
\item[(c)] Waarom is XRD beter geschikt dan spectroscopie voor kristallografische analyse?
\end{enumerate}

\item \textbf{Vacatures en temperatuur.} \\
De energie voor vacaturevorming in aluminium bedraagt $E_v = \SI{0.76}{\electronvolt}$.
\begin{enumerate}
\item[(a)] Bereken de evenwichtsvacaturefractie bij \SI{300}{\kelvin} en bij \SI{900}{\kelvin}.
\item[(b)] Verklaar waarom diffusie in metalen sterk toeneemt met temperatuur (verbind met Hoofdstuk~9).
\end{enumerate}

\item \textbf{Band gap en halfgeleiderclassificatie.} \\
Gegeven de band gaps:
\begin{itemize}
\item Si: \SI{1.12}{\electronvolt}
\item GaAs: \SI{1.42}{\electronvolt}
\item Diamant: \SI{5.5}{\electronvolt}
\item $\mathrm{TiO_2}$: \SI{3.2}{\electronvolt}
\end{itemize}
\begin{enumerate}
\item[(a)] Classificeer elk materiaal als geleider, halfgeleider of isolator.
\item[(b)] Bereken de minimale golflengte van licht dat $\mathrm{TiO_2}$ kan activeren. Verklaar waarom $\mathrm{TiO_2}$ werkt als fotokatalysator onder UV maar niet onder zichtbaar licht.
\item[(c)] Verbind met Hoofdstuk~2 (fotonenergie, $E = h\nu$) en Hoofdstuk~7 (katalyse).
\end{enumerate}

\item \textbf{Doping: n-type en p-type.} \\
Silicium wordt gedoteerd met fosfor (P) of boor (B).
\begin{enumerate}
\item[(a)] Welke doping levert n-type, welke p-type? Verklaar vanuit aantal valentie-elektronen.
\item[(b)] Bereken de ladingsdragerconcentratie bij \SI{1}{ppm} P in Si (veronderstel alle donoratomen ioniseren).
\item[(c)] Verklaar waarom gedoteerd Si bij kamertemperatuur geleidt, terwijl zuiver Si bij kamertemperatuur vrijwel niet.
\end{enumerate}

\item \textbf{Warmtecapaciteit: Dulong--Petit versus Debye.} \\
\begin{enumerate}
\item[(a)] Bereken $C_{V,m}$ van Pb bij \SI{300}{\kelvin} volgens Dulong--Petit. Vergelijk met experimentele waarde \SI{26.7}{\joule\per\mole\per\kelvin}.
\item[(b)] Verklaar waarom diamant bij \SI{298}{\kelvin} afwijkt ($C_{p,m} \approx \SI{6.1}{\joule\per\mole\per\kelvin}$).
\item[(c)] Bij welke temperatuur zal diamant $C_{V,m} \approx 3R$ benaderen?
\end{enumerate}

\item \textbf{Ge\"integreerde case: batterijmaterialen.} \\
$\mathrm{LiCoO_2}$ is een gelaagde structuur gebruikt in Li-ion batterijen.
\begin{enumerate}
\item[(a)] Verklaar waarom Li tussen de lagen kan intercaler zonder de structuur te breken (verbind met defecten).
\item[(b)] Bereken de theoretische capaciteit (in \si{Ah\per\kilogram}) als 1 Li per formule eenheid kan worden uitgewisseld (1 F per mol).
\item[(c)] Waarom worden moderne batterijen $\mathrm{LiFePO_4}$ (olivijnstructuur) gebruikt? (Verbind met kristalstructuur en stabiliteit.)
\item[(d)] Verbind met Hoofdstuk~6 (elektrochemie): de batterij is een herlaadbare galvanische cel.
\end{enumerate}

\end{enumerate}